% !TEX root = ../final-main.tex
\chapter{Technical Background}
\label{ch:technical-background}

\section{Overview of GaAs/AlGaAs HEMTs}

In GaAs/AlGaAs high-electron-mobility transistors (HEMTs), the interface
between the GaAs channel layer and the overlying AlGaAs barrier produces
a discontinuity in the conduction band due to the differing bandgaps of
the two materials. This offset creates a potential well on the GaAs side
of the junction, as illustrated in Figure~\ref{fig:band_bending}, which
confines electrons to the heterointerface and enables the formation of a
two-dimensional electron gas (2DEG).

\vspace{2mm}

Silicon donors in the AlGaAs layer ionise and release electrons, which
transfer across the heterointerface and accumulate in the potential
well on the GaAs side. This modulation-doping scheme leaves the
ionised donors in the AlGaAs while the mobile carriers reside in the
undoped GaAs channel, spatially separating the two. The resulting
suppression of ionised-impurity scattering yields the exceptionally
high electron mobility, high cut-off frequencies, low noise, and
strong transconductance that distinguish HEMTs from bulk-channel
field-effect transistors. \cite{liu1999}

\input{figures/fig-band-bending.tex}

\section{The GaAs/AlGaAs Heterostructure}

The devices investigated in this project employ a conventional
GaAs/AlGaAs heterostructure in which a sequence of epitaxial layers is
engineered to create and control a 2DEG. A schematic cross-section of
the structure is shown in Figure~\ref{fig:layer_structure}, highlighting
the cap, donor, spacer, and channel layers whose thicknesses and
compositions govern the electrostatic control and transport behaviour of
the device.

\input{figures/fig-layer-structure.tex}
\vspace{2mm}

\vspace{2mm}
At the surface of the structure lies a $10~\text{nm}$ thick, highly doped
GaAs cap layer. This cap provides a chemically stable surface for forming
the Schottky gate contact whilst protecting the underlying AlGaAs from oxidation. Its
high conductivity reduces series resistance in the access regions adjacent
to the source and drain contacts, improving carrier injection into the
channel.

\vspace{2mm}
Beneath the cap layer is a $40~\text{nm}$ thick silicon-doped AlGaAs
donor layer. Silicon donors act as electron suppliers, with carriers
transferring into the GaAs channel as a result of the conduction band
discontinuity at the heterointerface. This modulation-doping strategy
spatially separates the ionised dopants from the conducting channel.
Both the donor concentration and its thickness strongly influence the
resulting sheet carrier density of the 2DEG and therefore the threshold
voltage and maximum drain current of the device.

\vspace{2mm}
A $20~\text{nm}$ thick undoped AlGaAs spacer layer separates the donor
layer from the GaAs channel. This spacer suppresses Coulomb scattering
by physically isolating the 2DEG from ionised impurities, enabling
exceptionally high electron mobility. \cite{Hirakawa1986} The benefit is
particularly pronounced in HEMT devices operating at cryogenic
temperatures, where phonon scattering is reduced and impurity scattering
would otherwise dominate transport behaviour. \cite{Pfeiffer1989}
Increasing the spacer thickness improves mobility but weakens
electrostatic coupling between the gate and the channel; this trade-off
is a fundamental consideration in heterostructure design.

\vspace{2mm}
In this project, the thickness of the doped AlGaAs donor layer beneath
the gate represents the primary controllable parameter through which the
threshold voltage may be engineered.

\section{The Recessed-Gate Architecture}

In a typical HEMT device shown in Figure~\ref{fig:layer_structure}, the gate
is deposited directly onto the GaAs cap with the full donor layer intact beneath it.
The resulting threshold voltage is strongly negative, placing the device in depletion
mode. A method of realising enhancement-mode operation is to thin the donor layer by 
etching away the cap and penetrating a controlled depth into the AlGaAs, a process known as gate recession.

\vspace{2mm}

Figure~\ref{fig:recessed_gate} shows the resulting cross-section. The
etch opens a trench through the cap and into the donor layer, into which
the Schottky gate is deposited; the gate metal coats the walls of the trench and
rests on the cap surface either side of the trench, whilst the source and drain
contacts remain on the original surface, preserving low-resistance access
to the channel. The thickness of donor remaining beneath the gate, $d_d$,
is the critical process parameter: once it is reduced below a threshold
value, the zero-bias depletion region extends to the 2DEG and the device
switches to enhancement mode. 

% figures/fig-recessed-gate.tex
\begin{figure}[H]
    \vspace{2mm}
    \centering
    \begin{tikzpicture}[x=1.05cm, y=0.52cm, >=Stealth, font=\footnotesize]

        \def\w{8}
        \def\xrecL{2.4}   % trench left wall
        \def\xrecR{5.6}   % trench right wall
        \def\xgFL{2.2}    % gate flange left  (small overhang on cap)
        \def\xgFR{5.8}    % gate flange right

        % Bounding box
        \path[use as bounding box] (-2.2,-0.2) rectangle (\w+2.5,9.8);

        % Y-coordinates:
        %   0.0   substrate bottom
        %   1.5   channel bottom
        %   3.0   spacer bottom / 2DEG
        %   4.2   donor bottom
        %   5.95  trench floor  (d_d remaining beneath gate)
        %   6.7   cap bottom
        %   7.4   cap top / original surface
        %   7.9   gate flange top

        % --- Substrate ---
        \fill[gray!15] (0,0) rectangle (\w,1.5);
        \draw (0,0) rectangle (\w,1.5);
        \node at (4,0.75) {Semi-insulating GaAs substrate};

        % --- Channel ---
        \fill[blue!8] (0,1.5) rectangle (\w,3.0);
        \draw (0,1.5) rectangle (\w,3.0);
        \node at (4,2.25) {Undoped GaAs channel};

        % --- 2DEG ---
        \draw[blue,line width=1pt] (0,3.0) -- (\w,3.0);
        \node[blue,anchor=west] at (\w+0.15,3.0) {2DEG};

        % --- Spacer ---
        \fill[green!8] (0,3.0) rectangle (\w,4.2);
        \draw (0,3.0) rectangle (\w,4.2);
        \node at (4,3.6) {Undoped AlGaAs spacer};

        % --- Donor: single unified shape (U-profile) ---
        \fill[yellow!20]
            (0,4.2) -- (\w,4.2) -- (\w,6.7) --
            (\xrecR,6.7) -- (\xrecR,5.95) --
            (\xrecL,5.95) -- (\xrecL,6.7) --
            (0,6.7) -- cycle;
        \draw
            (0,4.2) -- (\w,4.2) -- (\w,6.7) --
            (\xrecR,6.7) -- (\xrecR,5.95) --
            (\xrecL,5.95) -- (\xrecL,6.7) --
            (0,6.7) -- cycle;
        \node[align=center] at (6.8,5.45) {Si-doped\\AlGaAs};

        % --- Cap: left of recess ---
        \fill[red!15] (0,6.7) rectangle (\xrecL,7.4);
        \draw (0,6.7) rectangle (\xrecL,7.4);

        % --- Cap: right of recess ---
        \fill[red!15] (\xrecR,6.7) rectangle (\w,7.4);
        \draw (\xrecR,6.7) rectangle (\w,7.4);
        \node[anchor=west] at (\w+0.15,7.05) {n$^{+}$ GaAs cap};

        % --- Gate metal: T-shape ---
        % Stem fills trench; small flange rests on cap surface
        \filldraw[fill=black!80, draw=black]
            (\xgFL,7.4) -- (\xgFL,7.9) --
            (\xgFR,7.9) -- (\xgFR,7.4) --
            (\xrecR,7.4) -- (\xrecR,5.95) --
            (\xrecL,5.95) -- (\xrecL,7.4) -- cycle;
        \node[anchor=south] at (4,7.9) {Gate};

        % --- Source ---
        \fill[black!30] (0.55,7.4) rectangle (1.25,7.95);
        \node[anchor=south] at (0.9,7.95) {Source};

        % --- Drain ---
        \fill[black!30] (6.75,7.4) rectangle (7.45,7.95);
        \node[anchor=south] at (7.1,7.95) {Drain};

        % --- Dimension: d_d ---
        \draw[<->] (\xrecL10,4.25) -- (\xrecL10,5.9);
        \node[anchor=west] at (\xrecL11,5.075) {$d_d$};

    \end{tikzpicture}
    \vspace{3mm}
    \caption{Cross-section of the recessed-gate GaAs/AlGaAs HEMT (not to
    scale).}
    \label{fig:recessed_gate}
\end{figure}


\vspace{2mm}

\section{Gate Control of Channel Charge and Threshold Voltage}

The electrical behaviour of the GaAs/AlGaAs HEMT is governed by the
field-effect interaction between the Schottky gate and the 2DEG. When a
metal gate is deposited onto the semiconductor surface, a depletion
region forms beneath it; the depth of this region is modulated by the
applied gate voltage ($V_G$). This modulation controls the sheet carrier
density ($n_s$) in the channel and, consequently, the drain current.
\cite{sze2021}

\subsection{Threshold Voltage and Required Recess Depth}
\label{sec:threshold-voltage}

The distinction between depletion-mode and enhancement-mode operation
is governed by the threshold voltage $V_{off}$. Within the analytical
framework of Lee \textit{et al.}~\cite{ivcv-hemt}:

\begin{equation}
    V_{off} = \phi_b - \Delta E_c - \frac{q N_d d_d^2}{2\varepsilon}
    \label{eq:threshold_voltage}
\end{equation}

\begin{quote}
    \footnotesize
    \textit{$\phi_b$: Schottky barrier height;
    $\Delta E_c$: conduction band offset;
    $N_d$: donor concentration;
    $d_d$: donor layer thickness;
    $\varepsilon$: permittivity of the AlGaAs barrier.}
\end{quote}

In the unrecessed device, $d_d = 40~\text{nm}$ and the quadratic term
in Equation~\ref{eq:threshold_voltage} dominates, yielding negative $V_{off}$ values;
the device is normally on and requires a negative gate bias to deplete the
channel. \cite{Mimura1980} Enhancement-mode operation requires $d_d$ to be
reduced until $V_{off}$ becomes positive.  

A critical donor layer thickness is derived by setting $V_{off} = 0$ and solving 
for $d_d$ in Equation~\ref{eq:threshold_voltage}.

\begin{equation}
    d_{d,\text{crit}} = \sqrt{\frac{2\varepsilon
    \left(\phi_b - \Delta E_c\right)}{q N_d}}
    \label{eq:crit_donor}
\end{equation}

\begin{table}[htbp]
    \centering
    \small
    \begin{tabular}{llll}
        \toprule
        Parameter               & Symbol           & Value                                    & Source              \\
        \midrule
        Schottky barrier height & $\phi_b$         & $0.80~\text{V}$                          & \cite{ivcv-hemt}    \\
        Conduction band offset  & $\Delta E_c$     & $0.23~\text{eV}$                         & \cite{liu1999}      \\
        Relative permittivity   & $\varepsilon_r$  & $12.2$                                   & \cite{liu1999}      \\
        Vacuum permittivity     & $\varepsilon_0$  & $8.854\times10^{-12}~\text{F\,m}^{-1}$  & Physical constant   \\
        Donor concentration     & $N_d$            & $1\times10^{18}~\text{cm}^{-3}$          & Wafer specification \\
        \bottomrule
    \end{tabular}
    \caption{Material parameters used to evaluate the critical donor
    thickness for the Al$_{0.3}$Ga$_{0.7}$As/GaAs heterostructure.}
    \label{tab:threshold_params}
\end{table}

\begin{equation}
    d_{d,\text{crit}} =
    \sqrt{\frac{2 \times 12.2 \times 8.854\times10^{-12} \times
    (0.80 - 0.23)}{1.6\times10^{-19} \times (1\times10^{18}) \times 10^{6}}}
    = \sqrt{7.70\times10^{-16}}
    \approx 28~\text{nm}
    \label{eq:crit_donor_eval}
\end{equation}

This is consistent with the $\sim\!25~\text{nm}$ quoted by Lee
\textit{et al.}~\cite{ivcv-hemt} for a similar heterostructure, with
the small difference attributable to the parameter values used.
The required total etch depth is:

\begin{equation}
    \text{Recess Depth} = d_{\text{cap}} +
    \left(d_{\text{donor}} - d_{d,\text{crit}}\right)
    = 10 + (40 - 28) = 22~\text{nm}
    \label{eq:recess_depth}
\end{equation}

The value of $22~\text{nm}$ serves as a theoretical reference rather
than a precise process target. In this project, multiple device sets are
fabricated with recess depths spanning approximately $10\text{--}30~\text{nm}$,
sweeping through the predicted threshold condition. Devices etched below
$22~\text{nm}$ are expected to exhibit progressively less negative $V_{off}$
as $d_d$ decreases; those etched beyond $22~\text{nm}$ should cross into
enhancement mode. This systematic variation allows the dependence of
threshold voltage on recess depth to be characterised
empirically, as well as any other affected characteristics such as
transconductance and drain current. As Equation~\ref{eq:crit_donor} is evaluated using literature
parameter values rather than direct measurement on these wafers, measurements of fabricated devices
will enable the validation of the theoretical prediction and
identify the true threshold depth for this specific heterostructure.

\subsection{Charge Distribution and Effective Gate Separation}
\label{sec:quantum-correction}

The threshold calculation above treats the 2DEG as a sheet of charge
located exactly at the GaAs/AlGaAs heterointerface. In the charge-control
model, the gate metal and the 2DEG are treated as the two plates of a
parallel-plate capacitor, separated by the total gate-to-channel distance
$d$. The capacitive coupling is assumed to occur through the depleted
AlGaAs barrier stack, with the gate voltage modulating the extent of the
depletion region and the model assigning a uniform permittivity
$\varepsilon$ to this region \cite{ivcv-hemt}.

\vspace{2mm}

In practice, however, the electrons forming the 2DEG occupy a quantised
triangular potential well on the GaAs side of the heterointerface. The
charge distribution therefore has finite thickness, and its centroid is
displaced into the GaAs rather than lying exactly at the interface. Lee
\textit{et al.} account for this by introducing an effective
gate-to-channel separation of $d + \Delta d$, where
$\Delta d \approx 80~\text{\AA}$ is a quantum-mechanical correction
associated with the spatial extent of the electron wavefunction
\cite{ivcv-hemt}. This correction enters the transconductance parameter
$\beta$ as

\begin{equation}
    \beta = \frac{\varepsilon \mu W}{(d + \Delta d)L}
    \label{eq:beta}
\end{equation}

\begin{quote}
    \footnotesize
    \textit{$\varepsilon$: permittivity of the barrier; $\mu$: electron
    mobility; $W$, $L$: gate width and length; $d$: physical
    gate-to-channel separation; $\Delta d$: effective charge-centroid
    offset.}
\end{quote}

\vspace{2mm}

The physical separation $d$ comprises the material lying between the gate
metal and the 2DEG. For the unrecessed device, this spans the full
epitaxial stack above the channel:
\begin{equation*}
    d_{\text{unrecessed}} = d_{\text{cap}} + d_{\text{donor}} +
    d_{\text{spacer}} = 10 + 40 + 20 = 70~\text{nm}
\end{equation*}

After recessing to the threshold condition
(Equation~\ref{eq:crit_donor_eval}), the cap is fully removed, the donor
layer is reduced to $d_{d,\text{crit}} = 28~\text{nm}$, and the spacer
remains intact:
\begin{equation*}
    d_{\text{recessed}} = d_{d,\text{crit}} + d_{\text{spacer}}
    = 28 + 20 = 48~\text{nm}
\end{equation*}

Including the charge-centroid correction therefore gives:

\begin{itemize}
    \item \textbf{Unrecessed gate:}
          $d_{\text{eff}} = 70 + 8 = 78~\text{nm}$, so the correction
          represents $8/78 \approx 10\%$ of the effective separation.
    \item \textbf{Recessed gate:}
          $d_{\text{eff}} = 48 + 8 = 56~\text{nm}$, so the correction
          represents $8/56 \approx 14\%$ of the effective separation.
\end{itemize}

The main relevance of $\Delta d$ in this project is that it modifies the
effective gate-to-channel separation entering the transconductance
parameter $\beta$. Physically, the correction reflects the fact that the
2DEG charge centroid lies inside the GaAs rather than exactly at the
heterointerface. Recessing the gate reduces the physical separation $d$,
but it does not alter the position of this charge centroid, so
$\Delta d$ becomes a larger fraction of the total effective separation as
the recess deepens. 

\vspace{2mm}

The correction is therefore relatively modest in the
unrecessed device, but becomes more significant once the gate is brought
closer to the channel. Although it does not change the basic
threshold-voltage argument developed previously, it is relevant when comparing
measured and expected transconductance, and more generally when
interpreting how recessing alters gate control of the channel.

\subsection{Gate Recess Risks and Limitations}

Gate recessing offers a practical route to threshold-voltage engineering, but it
also introduces process and device risks that must be carefully managed.

\vspace{2mm}

\noindent\textbf{Threshold-voltage sensitivity $(V_{off})$}

\vspace{2mm}

\noindent The quadratic dependence of $V_{off}$ on $d_d$
(Equation~\ref{eq:threshold_voltage}) means that small variations in recess
depth translate directly into shifts in threshold voltage. Near
$d_{d,\text{crit}}$, a $\pm 2~\text{nm}$ variation in etch depth can shift
$V_{off}$ by tens of millivolts \cite{ivcv-hemt}.

\vspace{2mm}

If an array of such devices were to be fabricated across an entire wafer for
subsequent circuit integration, the recess process would need to deliver
extremely high spatial uniformity in etch depth across the wafer, since even
small depth variations would produce threshold-voltage mismatch between
nominally identical transistors. An equivalent requirement would apply from
wafer to wafer if reproducible circuit behaviour is to be achieved in
manufacture. In logic circuits formed from multiple transistors, this
threshold-voltage spread can disrupt switching consistency, shift logic
thresholds, and reduce noise margins; in severe cases, the resulting mismatch
may substantially degrade circuit performance or prevent reliable operation
altogether.

\vspace{4mm}

\noindent \textbf{Increased Gate Leakage Current $(I_{g})$}

\noindent Gate recessing also introduces a greater risk of gate leakage current,
$I_{g}$. By depositing the gate metal into an etched trench, the gate is
brought into direct contact with the AlGaAs donor layer at an interface that
may have been roughened, oxidised, or otherwise chemically degraded during the
recess process. At the same time, reducing the effective gate-to-channel
separation increases the electric field across the barrier for a given applied
bias, which can promote tunnelling and elevate reverse-bias leakage currents.

$I_{g}$ is therefore influenced not only by the reduced physical
separation, but also by the quality of the recessed surface and the extent to
which the etch preserves a clean, well-controlled metal-semiconductor
interface.

\vspace{8mm}

Gate recessing is therefore most effective as a threshold-engineering tool when
both etch-depth control and interface quality are maintained to a high
standard. Quantifying this process sensitivity is an explicit objective of this
work. In particular, the repeatability of the $V_{off}$ and $I_{g}$ metrics
will be assessed across nominally identical devices on the same sample, as well
as across different samples fabricated at different times.
